\chapter{Discussão preliminar}

Com base no que foi disposto, verifica-se que a extração de sinais de fotopletismografia remota a partir de vídeos de pacientes internados em Unidades de Terapia Intensiva pode ser realizada por meio de métodos baseados em algoritmos e modelos matemáticos, denominados métodos não supervisionados, bem como pode ser realizada por meio de modelos baseados em aprendizado profundo, referidos como métodos supervisionados. Esses métodos resultam em uma onda de pulso de volume sanguíneo a qual possui  correspondência identificável com o sinal de fotopletismografia convencional e a partir da qual podem ser estimados valores de frequência cardíaca e de frequência respiratória. No entanto,  artefatos de iluminação e de movimento impactam de forma significativa na qualidade do sinal extraído, como foi possível verificar a partir dos resultados da Coleta Preliminar 2, na qual as estimativas de frequência cardíaca extraídas do vídeo com maior incidência desses elementos apresentaram, para a maioria dos métodos analisados, maior erro médio absoluto em relação aos valores de referência obtidos por meio do eletrocardiograma, quando comparado às estimativas advindas do vídeo em que o indivíduo monitorado permaneceu imóvel. De maneira complementar, a análise do erro médio por janela das estimativas de frequência cardíaca realizada tanto nos dados da Coleta Preliminar 1, quanto nos dados da Coleta Preliminar 2, permite identificar valores maiores de erro em janelas nas quais há mudanças expressivas de iluminação ou da posição do rosto, o que reforça o impacto desses artefatos. Além disso, artefatos de compressão de imagem também impactam negativamente no sinal de rPPG extraído, por isso a aquisição de dados foi realizada com a utilização de uma câmera industrial, com o intuito de minimizar a influência desse fator.

Ademais, o sinal que pode ser obtido por métodos de fotopletismografia remota é influenciado pela detecção e rastreamento da região de interesse, o que foi verificado no processamento dos dados da Coleta Preliminar 1. Nesse caso, o sinal extraído a partir do vídeo do paciente 1 apresentou um baixo índice de correlação com o sinal de fotopletismografia de referência, bem como as estimativas de frequência cardíaca realizadas resultaram nos maiores valores de erro médio absoluto entre todos os resultados. Isso pode ser explicado devido ao fato de que a ferramenta de detecção de face utilizada não reconheceu o rosto do paciente, similar a situações descritas por \citeonline{van2025camera} em relação ao monitoramento remoto em UTI, possivelmente devido à oclusão significativa dessa região, em decorrência da presença de equipamentos de suporte ventilatório. Assim, para todos os métodos, foram utilizados os quadros completos do vídeo como entrada, o que gera o aumento da incidência de elementos de ruído e, consequentemente, a degradação do sinal, o que é atenuado quando apenas a região de interesse é introduzida aos modelos. Observa-se, ainda, que nesse caso o sinal de PPG de referência difere da forma de onda fotopletismográfica convencional esperada, o que pode decorrer de fatores fisiológicos relacionados à condição clínica do paciente, de modo que são necessárias análises mais aprofundadas para esclarecer esse tipo de influência sobre o sinal obtido remotamente.

No que se refere às estimativas de frequência respiratória, nota-se que o sinal respiratório de referência do paciente 3 apresentou ruído de forma que o espectro de potência respectivo a esse sinal possui um espalhamento na banda de interesse e não possui um único pico distinguível (vide Apêndice \ref{apendice:spec-freqresp}). Nesse sentido, a ausência de um pico espectral dominante dificulta a definição inequívoca da frequência respiratória de referência, de modo que parte da discrepância observada entre as estimativas e a referência, nesse caso, pode estar associada à própria qualidade do sinal de referência, e não exclusivamente ao desempenho dos métodos avaliados.

Por fim, a partir dos valores de erro médio absoluto para as estimativas de frequência cardíaca e para as estimativas de frequência respiratória tanto a partir dos dados da Coleta Preliminar 1 quanto da Coleta Preliminar 2, nota-se que os métodos supervisionados apresentaram resultados comparáveis aos métodos não supervisionados, apesar de demandarem maior custo computacional para inferência. Ademais, entre os métodos supervisionados, modelos de mesma estrutura treinados em bases de dados diferentes resultaram em estimativas com erros expressivamente diferentes em relação à referência, o que demonstra a influência do conjunto de dados de treinamento no desempenho dos modelos em situações realistas. Por conseguinte, ambos os tipos de método de extração de sinais de fotopletismografia remota apresentam métricas de erro e valores para a razão sinal-ruído que reafirmam o potencial de aplicação desse tipo de tecnologia em situações clínicas. Faz-se necessário, todavia, a análise em uma quantidade maior de amostras para que conclusões mais detalhadas possam ser esclarecidas. 

